\chapter{Operadores}

\section{Aritméticos}

\begin{center}
\begin{tabular}{cll}
\toprule
\textbf{Operador} & \textbf{Operação} & \textbf{Exemplo} \\
\midrule
\texttt{+}  & Adição               & \texttt{3 + 2} → \texttt{5} \\
\texttt{-}  & Subtração            & \texttt{5 - 3} → \texttt{2} \\
\texttt{*}  & Multiplicação        & \texttt{4 * 3} → \texttt{12} \\
\texttt{/}  & Divisão              & \texttt{7 / 2} → \texttt{3} (inteiros) \\
\texttt{\^} & Potência             & \texttt{2\^{}10} → \texttt{1024} \\
\texttt{\%} & Módulo               & \texttt{7 \% 3} → \texttt{1} \\
\bottomrule
\end{tabular}
\end{center}

Divisão entre dois \texttt{Int} é inteira (truncamento em direção a zero). Para
divisão real, ao menos um operando deve ser \texttt{Float}:

\begin{lstlisting}
7 / 2      // 3
7 / 2.0    // 3.5
7.0 / 2    // 3.5
\end{lstlisting}

\section{Comparação}

\begin{center}
\begin{tabular}{cll}
\toprule
\textbf{Operador} & \textbf{Significado} & \textbf{Exemplo} \\
\midrule
\texttt{==} & Igual               & \texttt{3 == 3} → \texttt{true} \\
\texttt{!=} & Diferente           & \texttt{3 != 4} → \texttt{true} \\
\texttt{>}  & Maior               & \texttt{5 > 3}  → \texttt{true} \\
\texttt{<}  & Menor               & \texttt{3 < 5}  → \texttt{true} \\
\texttt{>=} & Maior ou igual      & \texttt{3 >= 3} → \texttt{true} \\
\texttt{<=} & Menor ou igual      & \texttt{2 <= 3} → \texttt{true} \\
\bottomrule
\end{tabular}
\end{center}

\section{Lógicos}

\begin{lstlisting}
true and false    // false
true or false     // true
not true          // false
\end{lstlisting}

\lstinline{and} e \lstinline{or} têm avaliação em curto-circuito: o segundo
operando não é avaliado se o resultado já é determinado pelo primeiro.

\section{Operador \texttt{in}}

Verifica pertencimento a uma lista:

\begin{lstlisting}
let regioes = ["norte", "sul", "centro"]
"sul" in regioes      // true
"leste" in regioes    // false
\end{lstlisting}

Funciona também com ranges:

\begin{lstlisting}
5 in 1..10    // true
15 in 1..10   // false
\end{lstlisting}

\section{Operador pipe \texttt{|>}}
\label{sec:pipe}

O operador \lstinline{|>} passa o valor da esquerda como primeiro argumento da
função à direita:

\begin{lstlisting}
// sem pipe
display filter(df, x > 5)

// com pipe — mais legível
df |> filter(x > 5) |> display
\end{lstlisting}

O pipe tem semântica especial com DataFrames: quando usado sem atribuição explícita
(\lstinline{let}), o resultado é atribuído de volta à variável de origem:

\begin{lstlisting}
let df = load("dados.csv")

// atualiza df no lugar
df |> filter(x > 0) |> sort(x)

// cria df2, df permanece inalterado
let df2 = df |> filter(x > 0)
\end{lstlisting}

Este comportamento é discutido em detalhe no Capítulo~\ref{cap:pipes}.

\section{Precedência}

Da menor para a maior prioridade:

\begin{enumerate}
  \item \texttt{or}
  \item \texttt{and}
  \item \texttt{not}
  \item \texttt{== != > < >= <=}
  \item \texttt{+ -}
  \item \texttt{* / \%}
  \item \texttt{\^} (associa à direita)
  \item Unário \texttt{-}
  \item Chamada de função, indexação, acesso a campo
\end{enumerate}

Use parênteses para clareza quando a precedência não for óbvia:

\begin{lstlisting}
2 + 3 * 4       // 14  (multiplicação primeiro)
(2 + 3) * 4     // 20
\end{lstlisting}
